% =====================================================================
\chapter{Introduction}
\label{ch:intro}
% =====================================================================

The spread of generative artificial intelligence has changed how written
content is produced. Large language models (LLMs) built on the transformer
architecture~\cite{vaswani2017attention} and trained on internet-scale corpora
can now produce text that is fluent, coherent and, in many settings,
stylistically indistinguishable from human writing~\cite{radford2019language,
openai2023gpt4,touvron2023llama}. This capability is valuable for productivity
and creativity, but it also creates serious problems for academic integrity,
for the verification of originality, and for the trustworthiness of digital
information~\cite{tang2024science,wu2025survey}.

This chapter introduces the problem addressed by the project, the motivation
behind it, and the theoretical concepts on which the proposed system,
\textbf{\veriscript}, is built. It closes with the scope of the work and the
organisation of the rest of the report.

% ---------------------------------------------------------------------
\section{Background}
% ---------------------------------------------------------------------

\subsection{The Rise of Large Language Models}

Statistical language modelling has a long history in natural language
processing, but its capability changed sharply with the introduction of
self-attention~\cite{vaswani2017attention}. GPT-2~\cite{radford2019language}
showed that a decoder-only transformer trained on a large web corpus could
produce paragraphs of convincing prose with no task-specific supervision.
Successive generations of models, including GPT-4~\cite{openai2023gpt4}, the
open-weight LLaMA family~\cite{touvron2023llama}, and commercial assistants
such as Claude, have scaled the number of parameters, the size of the training
data and the use of instruction tuning. The result is a class of systems that
write essays, answer examination questions, summarise articles and imitate a
requested style on demand.

These systems are now inexpensive and widely available. A student can obtain
a complete assignment in seconds through a free chat interface, and a content
farm can generate thousands of articles per day through an application
programming interface (API). The cost of producing fluent text has therefore
fallen almost to zero, while the cost of verifying who wrote a given piece of
text has not.

\subsection{Impact on Academic Integrity and the Information Ecosystem}

Educational institutions rely on written work as evidence of a student's
understanding. When an assignment can be produced by a model, written
submissions lose much of their value as an assessment instrument unless the
institution has some means of checking authorship. Publishers, peer-review
systems, recruiters and online platforms face a similar problem: reviews,
news articles, cover letters and social media posts can all be produced at
scale, and readers can no longer assume that fluent text reflects human effort
or human accountability~\cite{zellers2019grover,tang2024science}.

Because of this, AI-generated text detection has become a distinct area of
research and a commercial market. Institutions want a tool that answers a
simple question, ``was this text written by a person or by a model?'', quickly,
cheaply, accurately and in a way that can be explained to the person whose
work is being judged.

\subsection{Why Conventional Plagiarism Detection Is Not Enough}

Traditional plagiarism detectors compare a submitted document against a
reference database of existing documents and report passages that overlap.
This works when the dishonest submission is copied from somewhere. AI-generated
text, however, is sampled afresh from a probability distribution every time a
model is prompted: it is normally \emph{novel} and does not exist in any prior
corpus. Overlap-based detection is therefore structurally unable to address
the problem. Detectors of machine-generated text must instead rely on
statistical and learned \emph{signatures} of the generation process itself
rather than on textual overlap with known sources~\cite{wu2025survey}.

% ---------------------------------------------------------------------
\section{Motivation and Perspective}
% ---------------------------------------------------------------------

Two broad families of detectors have emerged in the literature and in
commercial products. The first family is \textbf{heuristic or statistical
detection}. Because an LLM samples each token from a learned probability
distribution, its output tends to have lower perplexity and lower
\emph{burstiness} (more uniform variation from sentence to sentence) than
human writing~\cite{gehrmann2019gltr,mitchell2023detectgpt,hans2024binoculars}.
Detectors built on these signals are fast, need little or no training data and
are fully interpretable, but their accuracy is limited and they are sensitive
to changes of domain and to adversarial paraphrasing~\cite{krishna2023paraphrasing}.

The second family is \textbf{transformer-based classification}, in which a
pretrained encoder such as RoBERTa~\cite{liu2019roberta} or
DeBERTa~\cite{he2021deberta,he2023debertav3} is fine-tuned as a binary
classifier on labelled human and machine text~\cite{solaiman2019release,
guo2023hc3}. Such classifiers reach very high in-domain accuracy, but every
prediction needs a full forward pass through a network with over a hundred
million parameters. For an institution that must screen thousands of
submissions at the end of a term, often without a dedicated graphics processing
unit (GPU), this cost is significant. Commercial services built on such models
add licence fees and per-page charges, and most of them report only a score
without explaining how it was obtained.

Institutions and platforms are therefore forced to choose between a detector
that is fast and cheap but shallow, and one that is accurate but slow,
expensive and opaque. Two further problems make this choice harder. First,
large benchmarks show that most detectors generalise poorly to generators that
were not present in their training data~\cite{dugan2024raid,wang2024m4,
li2024mage}. Second, widely used detectors have been shown to misclassify the
writing of non-native English speakers as AI-generated far more often than
native writing~\cite{liang2023biased}, which is a serious fairness risk when
the outcome is an accusation of academic misconduct.

The perspective taken in this project is that the two families are not rivals
but \emph{complements}. Most documents are easy: a heuristic model can label
them correctly with high confidence. Only a minority are genuinely ambiguous,
and only those need the expensive, context-aware judgement of a transformer.
A system that spends its computational budget where it is needed can,
in principle, approach the accuracy of the transformer at close to the cost of
the heuristic. This is the motivation for \veriscript.

% ---------------------------------------------------------------------
\section{Description of Theoretical Concepts}
\label{sec:theory}
% ---------------------------------------------------------------------

This section introduces the concepts used throughout the report: how
language models generate text, the statistics that distinguish model output
from human writing, transformer classifiers, probability calibration and
classifier cascades.

\subsection{Autoregressive Language Modelling}

A document $d$ is a sequence of tokens $x_1, x_2, \ldots, x_N$. An
autoregressive language model with parameters $\theta$ factorises the
probability of the sequence as a product of next-token conditionals:
\begin{equation}
  P_\theta(d) \;=\; \prod_{i=1}^{N} P_\theta\!\left(x_i \mid x_1,\ldots,x_{i-1}\right).
  \label{eq:autoregressive}
\end{equation}
During generation the model repeatedly computes the distribution over the
vocabulary for the next position, selects a token, appends it to the context
and repeats, as shown in Fig.~\ref{fig:autoregressive}. Training adjusts
$\theta$ to maximise the likelihood of a very large corpus of human text, so
the model learns which continuations are \emph{typical}.

\begin{figure}[H]
\centering
\begin{tikzpicture}[node distance=0.9cm and 1.0cm]
  \node[blk,minimum width=3.4cm] (ctx) {Context\\$x_1,\ldots,x_{i-1}$};
  \node[blk,right=of ctx,minimum width=3.0cm] (lm) {Transformer\\decoder ($\theta$)};
  \node[blk,right=of lm,minimum width=3.0cm] (dist) {Distribution\\$P_\theta(\,\cdot \mid x_{<i})$};
  \node[blk,below=1.2cm of dist,minimum width=3.0cm] (samp) {Decoding rule\\(greedy, top-$k$, top-$p$)};
  \node[blk,below=1.2cm of ctx,minimum width=3.4cm] (tok) {Selected token $x_i$};
  \draw[arr] (ctx) -- (lm);
  \draw[arr] (lm) -- (dist);
  \draw[arr] (dist) -- (samp);
  \draw[arr] (samp) -- (tok);
  \draw[arr] (tok) -- node[lbl,left] {append} (ctx);
\end{tikzpicture}
\caption{Autoregressive generation loop of a large language model}
\label{fig:autoregressive}
\end{figure}

\subsection{Decoding Strategies}

The way the next token is chosen from $P_\theta(\cdot \mid x_{<i})$ strongly
affects the statistical character of generated text. \emph{Greedy decoding}
always picks the most probable token and tends to produce repetitive text.
\emph{Temperature sampling} rescales the logits $z_v$ before the softmax,
\begin{equation}
  P_T(v \mid x_{<i}) = \frac{\exp(z_v/T)}{\sum_{u}\exp(z_u/T)},
  \label{eq:temperature}
\end{equation}
so that $T<1$ sharpens and $T>1$ flattens the distribution. \emph{Top-$k$}
sampling restricts the choice to the $k$ most likely tokens, and \emph{nucleus}
(top-$p$) sampling restricts it to the smallest set of tokens whose cumulative
probability exceeds $p$~\cite{holtzman2020curious}. All of these rules remove
the low-probability tail of the distribution. Human writers do not obey such a
rule; they regularly choose words that a model would consider unlikely. This
asymmetry is the root of most statistical detection signals~\cite{gehrmann2019gltr,
ippolito2020automatic}.

\subsection{Perplexity and Token-Level Statistics}

Given a reference language model, the \emph{perplexity} (PPL) of a document
is the exponential of the average negative log-likelihood of its tokens:
\begin{equation}
  \mathrm{PPL}(d) \;=\; \exp\!\left(-\frac{1}{N}\sum_{i=1}^{N}
      \log P_\theta\!\left(x_i \mid x_{<i}\right)\right).
  \label{eq:perplexity}
\end{equation}
Low perplexity means that the reference model found the text predictable.
Because generated text is assembled from high-probability choices, it usually
has lower perplexity than human text on the same topic. Related token-level
statistics include the \emph{rank} of each observed token in the predicted
distribution and the \emph{entropy} of the predicted distribution,
\begin{equation}
  H_i \;=\; -\sum_{v \in \mathcal{V}} P_\theta(v \mid x_{<i})\,\log P_\theta(v \mid x_{<i}),
  \label{eq:entropy}
\end{equation}
both of which were used by GLTR~\cite{gehrmann2019gltr}. Perplexity on its own
is not a reliable detector, because predictable human writing (formal,
templated or simplified prose) also has low perplexity~\cite{liang2023biased}.
It is, however, a cheap and informative \emph{feature}.

\subsection{Burstiness and Stylometric Features}

Human writing is \emph{bursty}: a long, complex sentence is often followed by a
short one, and the predictability of the text rises and falls across a
document. Model output tends to be more uniform. Let a document contain $S$
sentences of lengths $\ell_1,\ldots,\ell_S$ with mean $\mu_\ell$ and standard
deviation $\sigma_\ell$. A common burstiness coefficient is
\begin{equation}
  \beta \;=\; \frac{\sigma_\ell - \mu_\ell}{\sigma_\ell + \mu_\ell},
  \qquad -1 \le \beta \le 1,
  \label{eq:burstiness}
\end{equation}
where values near $-1$ indicate very regular sentence lengths and larger values
indicate bursty text. The same idea can be applied to the per-sentence
perplexity, giving a measure of how much the predictability fluctuates.

\emph{Stylometric} features describe the style of a text independently of its
topic. Examples are the type-token ratio
\begin{equation}
  \mathrm{TTR}(d) = \frac{V}{N},
  \label{eq:ttr}
\end{equation}
where $V$ is the number of distinct word types, the distribution of
part-of-speech (POS) $n$-grams, average word length, punctuation frequency and
the rate of function words. These features are cheap to compute and easy to
explain, which makes them suitable for the first stage of a detection
pipeline.

\subsection{Transformer Encoders and Fine-Tuned Classification}

A transformer encoder maps a token sequence to a sequence of contextual
vectors using stacked self-attention layers~\cite{vaswani2017attention}. For
query, key and value matrices $Q$, $K$ and $V$ of dimension $d_k$, one
attention head computes
\begin{equation}
  \mathrm{Attention}(Q,K,V) \;=\; \mathrm{softmax}\!\left(\frac{QK^{\top}}{\sqrt{d_k}}\right)V.
  \label{eq:attention}
\end{equation}
BERT~\cite{devlin2019bert} pretrains such an encoder with a masked-language
objective; RoBERTa~\cite{liu2019roberta} improves the pretraining recipe; and
DeBERTa~\cite{he2021deberta} separates content and position information in the
attention computation, with DeBERTa-v3~\cite{he2023debertav3} adding a more
efficient pretraining objective. For classification, a special token
(\texttt{[CLS]} or \texttt{<s>}) is placed at the start of the input and its
final hidden vector $\mathbf{h}_{\mathrm{CLS}}$ is passed to a small output
layer, as in Fig.~\ref{fig:encoder}:
\begin{equation}
  p_t(d) \;=\; \sigma\!\left(\mathbf{w}_t^{\top}\mathbf{h}_{\mathrm{CLS}} + b_t\right).
  \label{eq:clshead}
\end{equation}
The whole network is then fine-tuned on labelled examples with the
cross-entropy loss. Because the encoder reads the entire context in both
directions, it can learn subtle cues of discourse structure and phrasing that
hand-crafted statistics miss.

\begin{figure}[H]
\centering
\begin{tikzpicture}[node distance=0.55cm]
  \node[blk,minimum width=10.5cm] (tok) {Tokenizer: \texttt{<s>} $x_1\;x_2\;\cdots\;x_L$ \texttt{</s>}};
  \node[blk,below=of tok,minimum width=10.5cm] (emb) {Token + position embeddings};
  \node[blk,below=of emb,minimum width=10.5cm,minimum height=1.4cm] (enc)
     {Stacked self-attention encoder layers (12 layers for the base models)};
  \node[blk,below=of enc,minimum width=4.2cm] (cls) {$\mathbf{h}_{\mathrm{CLS}}$ (pooled vector)};
  \node[blk,below=of cls,minimum width=4.2cm] (head) {Linear layer + sigmoid};
  \node[term,below=of head] (out) {$p_t(d)$ = P(AI-generated)};
  \draw[arr] (tok) -- (emb);
  \draw[arr] (emb) -- (enc);
  \draw[arr] (enc) -- (cls);
  \draw[arr] (cls) -- (head);
  \draw[arr] (head) -- (out);
\end{tikzpicture}
\caption{Fine-tuned transformer encoder used as a binary text classifier}
\label{fig:encoder}
\end{figure}

\subsection{Probability Calibration}

A classifier is \emph{calibrated} if, among all documents to which it assigns
probability $0.8$, about 80\% are in fact AI-generated. Calibration matters for
this project because the cascade compares the heuristic stage's probability
against fixed thresholds; if that probability is badly calibrated, the
thresholds do not mean what they appear to mean. Modern neural networks are
often over-confident~\cite{guo2017calibration}. A standard correction is
\emph{Platt scaling}~\cite{platt1999probabilistic}, which fits two scalars
$a$ and $b$ on a held-out set so that the raw score $s(d)$ becomes
\begin{equation}
  p(d) \;=\; \sigma\!\left(a\,s(d) + b\right).
  \label{eq:platt}
\end{equation}
Calibration quality is commonly measured by the expected calibration error
(ECE), which bins predictions by confidence and averages the gap between
confidence and observed accuracy across bins~\cite{guo2017calibration}.

\subsection{Classifier Cascades and Early Exit}

A \emph{cascade} is a sequence of classifiers of increasing cost in which each
input is passed to the next stage only if the current stage is not confident
enough. The idea is well established in computer vision, where the
Viola--Jones face detector rejects most image windows with a few cheap features
and spends computation only on promising regions~\cite{viola2001rapid}. In deep
learning the same principle appears as \emph{early exit}: BranchyNet adds
side classifiers to intermediate layers of a network~\cite{teerapittayanon2016branchynet},
and DeeBERT applies dynamic early exiting to BERT~\cite{xin2020deebert}. More
recently, FrugalGPT showed that cascading LLM APIs of increasing price can
match the best model's quality at a fraction of its cost~\cite{chen2023frugalgpt}.

For a two-stage cascade with per-document costs $C_h$ (cheap stage) and $C_t$
(expensive stage), and an escalation rate $\alpha$ (the fraction of documents
passed to the second stage), the expected cost per document is
\begin{equation}
  \mathbb{E}[C] \;=\; C_h + \alpha\, C_t .
  \label{eq:cascadecost}
\end{equation}
When $C_t \gg C_h$ and $\alpha$ is small, the cascade costs only a little more
than the cheap stage alone. The accuracy of the cascade depends on whether the
cheap stage is right on the documents it keeps, which in turn depends on its
calibration. Fig.~\ref{fig:cascadeconcept} illustrates the idea in the setting
of this project.

\begin{figure}[H]
\centering
\begin{tikzpicture}[node distance=0.8cm and 1.1cm]
  \node[io,minimum width=2.6cm] (in) {All documents};
  \node[blk,right=of in,minimum width=3.0cm] (s1) {Stage 1\\cheap, cost $C_h$};
  \node[blk,right=1.6cm of s1,minimum width=3.0cm] (s2) {Stage 2\\expensive, cost $C_t$};
  \node[term,below=1.1cm of s1] (o1) {Confident: decide now\\(fraction $1-\alpha$)};
  \node[term,below=1.1cm of s2] (o2) {Final decision\\(fraction $\alpha$)};
  \draw[arr] (in) -- (s1);
  \draw[arr] (s1) -- node[lbl,above] {uncertain} (s2);
  \draw[arr] (s1) -- (o1);
  \draw[arr] (s2) -- (o2);
\end{tikzpicture}
\caption{General principle of a two-stage classifier cascade}
\label{fig:cascadeconcept}
\end{figure}

% ---------------------------------------------------------------------
\section{Families of AI-Generated Text Detection}
% ---------------------------------------------------------------------

The detection methods discussed in this report fall into four families,
summarised in Fig.~\ref{fig:taxonomy}: zero-shot statistical methods that use
a language model's own probabilities without training a classifier;
feature-based supervised methods that train a light classifier on hand-crafted
or model-derived features; fine-tuned neural classifiers; and watermarking,
which requires cooperation from the model provider at generation time.
Table~\ref{tab:families} compares the families on the properties that matter
for deployment.

\begin{figure}[H]
\centering
\begin{tikzpicture}[
   level distance=1.7cm,
   level 1/.style={sibling distance=3.7cm},
   level 2/.style={sibling distance=1.9cm},
   every node/.style={blk,font=\footnotesize,minimum height=0.8cm,text width=2.9cm},
   edge from parent/.style={draw,thick}]
\node[text width=4.4cm] {AI-generated text detection}
  child {node {Zero-shot statistical}
     child {node[text width=1.6cm] {GLTR}}
     child {node[text width=1.6cm] {DetectGPT, Binoculars}}}
  child {node {Feature-based supervised}
     child {node[text width=1.6cm] {Stylometry + ML}}
     child {node[text width=1.6cm] {Ghost-buster}}}
  child {node {Fine-tuned neural classifiers}
     child {node[text width=1.6cm] {RoBERTa}}
     child {node[text width=1.6cm] {DeBERTa, RADAR}}}
  child {node {Generation-time watermarking}
     child {node[text width=1.6cm] {Green/red token lists}}};
\end{tikzpicture}
\caption{Taxonomy of AI-generated text detection approaches}
\label{fig:taxonomy}
\end{figure}

\begin{table}[H]
\caption{Comparison of detection families on deployment-relevant properties}
\label{tab:families}
\centering
\small
\begin{tabularx}{\textwidth}{|p{2.9cm}|C|C|C|C|}
\hline
\textbf{Property} & \textbf{Zero-shot statistical} & \textbf{Feature-based supervised}
  & \textbf{Fine-tuned transformer} & \textbf{Watermarking} \\ \hline
Training data needed & None & Small & Large labelled corpus & None (at detection) \\ \hline
Inference cost & Low to high$^{\ast}$ & Low & High & Very low \\ \hline
In-domain accuracy & Moderate to high & Moderate & Very high & Very high \\ \hline
Interpretability & High & High & Low & High \\ \hline
Robustness to new generators & Moderate & Moderate & Low to moderate & Provider-specific \\ \hline
Works on arbitrary text & Yes & Yes & Yes & No \\ \hline
\multicolumn{5}{|l|}{\footnotesize $^{\ast}$DetectGPT needs many perturbations per document; GLTR-style statistics need one pass.} \\ \hline
\end{tabularx}
\end{table}

The table shows that no single family is best on every property. Feature-based
methods are cheap and interpretable; fine-tuned transformers are accurate.
\veriscript{} combines the two in a cascade so that each is used for the
documents on which it is most effective.

% ---------------------------------------------------------------------
\section{Challenges in Detecting AI-Generated Text}
% ---------------------------------------------------------------------

Several properties of the problem make detection harder than an ordinary
text-classification task. They shape the design decisions made later in this
report.

\subsection{A Shrinking Distributional Gap}

Each new generation of language models is trained to imitate human text more
closely. As the model distribution approaches the human distribution, the
statistical differences that detectors rely on become smaller, and in the limit
no detector can do better than chance~\cite{sadasivan2023reliably}. In practice
current models still leave measurable traces, but a detector must be
re-evaluated as new generators appear, which is why this project evaluates
generalisation to a held-out generator explicitly.

\subsection{Editing and Paraphrasing}

Generated text is rarely submitted unchanged. A student may edit the output,
merge it with their own writing, or pass it through a paraphrasing tool.
Paraphrasing models designed for this purpose can reduce the accuracy of
strong detectors dramatically while preserving meaning~\cite{krishna2023paraphrasing}.
Lightly edited text tends to move the heuristic statistics towards the human
range, which is one reason why ambiguous cases should be referred to a model
that reads the whole context.

\subsection{Short Texts and Domain Shift}

Statistical evidence accumulates with length. A paragraph of 50 words contains
too few tokens and sentences for perplexity or burstiness to be estimated
reliably, so detectors are much less accurate on short texts. Detectors also
behave differently across domains: technical answers, legal prose and poetry
have very different baseline statistics, and a detector calibrated on one
domain can be badly mis-calibrated on another~\cite{wang2024m4,li2024mage}.

\subsection{The Base-Rate Problem}

Even an accurate detector produces many false accusations if AI-generated
submissions are rare. Let $\pi$ be the fraction of submissions that are
AI-generated, and let the detector have true positive rate TPR and false
positive rate FPR. By Bayes' theorem the probability that a flagged
submission really is AI-generated (the precision) is
\begin{equation}
  P(\text{AI} \mid \text{flagged}) = \frac{\mathrm{TPR}\cdot\pi}{\mathrm{TPR}\cdot\pi + \mathrm{FPR}\cdot(1-\pi)} .
  \label{eq:baserate}
\end{equation}
For example, with $\mathrm{TPR}=0.95$, $\mathrm{FPR}=0.05$ and $\pi = 0.10$,
the precision is $0.095/(0.095+0.045) \approx 0.68$: roughly one flagged
submission in three would be a false accusation, even though the detector is
``95\% accurate''. The false positive rate therefore matters far more than
headline accuracy in academic use, and it must be measured separately for
groups of writers who may be affected differently~\cite{liang2023biased}.
This is why \veriscript{} treats its output as a calibrated probability to be
reviewed by a person, and why fairness is a primary evaluation criterion.

\subsection{Cost and Transparency at Scale}

Finally, a detector is useful only if it can be run on every submission and if
its decisions can be explained. Accurate neural detectors are expensive to run
over large batches, and commercial services are closed. The combination of
these practical constraints with the statistical difficulties above is the
setting in which \veriscript{} is designed.

% ---------------------------------------------------------------------
\section{Applications}
% ---------------------------------------------------------------------

An open, inexpensive and explainable detector is useful in several settings:
\begin{itemize}
  \item \textbf{Academic institutions}: screening assignments, theses and
        examination answers in bulk, with explanations that can be shared
        with the student during a review;
  \item \textbf{Admissions and recruitment}: triaging statements of purpose,
        cover letters and essays before human reading;
  \item \textbf{Publishing and peer review}: checking submissions and reviews
        for undisclosed machine generation;
  \item \textbf{Online platforms}: flagging bulk-generated reviews, comments
        and articles for moderator attention;
  \item \textbf{Research}: providing an open, reproducible baseline against
        which new detection methods can be compared on both accuracy and cost.
\end{itemize}
In every case the detector acts as a first filter that directs human
attention; it is not a substitute for human judgement.

% ---------------------------------------------------------------------
\section{Scope of the Project}
% ---------------------------------------------------------------------

The project designs, implements and evaluates a hybrid detector for English
prose. The following are \textbf{within scope}:
\begin{itemize}
  \item binary classification of a document as human-written or AI-generated,
        with a calibrated probability and a short explanation;
  \item a heuristic detector based on perplexity, burstiness and stylometric
        features with a classical machine-learning head;
  \item a fine-tuned RoBERTa-base classifier (the code also supports DeBERTa-v3
        as a drop-in alternative);
  \item a confidence-gated cascade whose ambiguity band is tuned on validation
        data;
  \item evaluation of accuracy, efficiency, cross-generator generalisation,
        robustness and fairness on public benchmarks (HC3, M4, RAID, MAGE and
        the Liang \textit{et al.} essay sets) covering eight generator
        families, including GPT-4, Llama and Mistral models;
  \item an open-source command line interface (CLI), REST service and web
        demonstration, together with a pipeline for adding custom generations
        from GPT-4, Claude and Llama models.
\end{itemize}

The following are \textbf{outside the scope} of this project: languages other than
English, mixed documents in which only some sentences are machine-generated
(sentence-level localisation), source attribution (identifying \emph{which}
model wrote a text), watermark detection, and any use of the system as the sole
basis for disciplinary decisions. The system is intended to support a human
reviewer, not to replace one.

% ---------------------------------------------------------------------
\section{Organisation of the Report}
% ---------------------------------------------------------------------

The rest of this report is organised as follows.
\textbf{Chapter~2} reviews related work in six areas: statistical and heuristic
methods, transformer-based classifiers, commercial tools, watermarking,
benchmark datasets and studies of reliability and fairness, and identifies the
research gap. \textbf{Chapter~3} formulates the problem, states it precisely,
depicts its inputs, processes and outputs, and lists the objectives.
\textbf{Chapter~4} presents the proposed methodology, the cascade algorithm and
the threshold-tuning procedure. \textbf{Chapter~5} gives the system design
through data flow diagrams and UML models. \textbf{Chapter~6} describes the
implementation: the algorithms, software tools, experimental setup, the
corpus that was built and the trained models. \textbf{Chapter~7} analyses the
results: accuracy, calibration, robustness, generalisation, fairness and
efficiency of the three systems, measured against the success criteria and
hypotheses of Chapter~3. \textbf{Chapter~8} draws conclusions, states the
limitations of the work and outlines future scope. The appendices contain the
contribution of the project and a short user guide.
